La segunda relación fundamental aparece cuando se comparan acción y gravedad.\par

En HMT, la acción mide cuánto ocurre y conserva simultáneamente la orientación de la realización. Hay una sección anterior al retorno y otra posterior. La fase puede volver acompañada por una historia de memoria diferente. Esa diferencia queda almacenada en la razón bidireccional de las dos secciones de acción.\par

La gravedad responde de manera recíproca.\par

El transductor gravitatorio tiene la forma\par

\[
G_a(\theta)
=
\theta\frac{c_{\rm int}^{5}t_0^2}{\hbar_a}.
\]

La acción aparece en el denominador. Por tanto, si al cambiar de hoja \(\hbar\) se multiplica por una razón \(R_{\rm act}\), la constante gravitatoria se multiplica por la razón inversa:\par

\[
\frac{\hbar_{\rm pre}}{\hbar_{\rm ret}}
=
R_{\rm act},
\qquad
\frac{G_{\rm pre}}{G_{\rm ret}}
=
R_{\rm act}^{-1}.
\]

La acción y la gravedad forman aquí una pareja contragrediente: cuando una avanza en una dirección, la otra compensa exactamente ese desplazamiento.\par

Lo que permanece invariante es\par

\[
G_a\hbar_a.
\]

Y ésa es precisamente la combinación que determina el área de Planck:\par

\[
\ell_P^2=\frac{G_a\hbar_a}{c_{\rm int}^3}.
\]

La acción recuerda la hoja. La gravedad responde a esa memoria. El área permanece invariante.\par

Ésta es una idea extraordinariamente fértil. El área de Planck emerge como el objeto que sobrevive cuando acción y gravedad se transforman recíprocamente.\par

Podría decirse así: \(\hbar\) conserva la orientación dinámica; \(G\) conserva la compensación geométrica; \(\ell_P^2\) conserva aquello en lo que ambas realizaciones están de acuerdo.\par

El área es, por tanto, una moneda de cambio entre física cuántica y geometría, porque la transformación bidireccional de \(\hbar\) y \(G\) selecciona exactamente esa combinación como invariante.\par

Esta reciprocidad permite comprender de otra forma la relación entre información y gravedad. La información necesita una distinción de estados; la acción mide el coste de transición entre ellos. Pero si esa transición debe adquirir una geometría, la respuesta gravitatoria cambia de manera inversa para mantener estable el cuanto de área.\par

La geometría incorpora la información de fase mediante una compensación exacta.\par

Desde una lectura machiana, éste es ya un resultado muy sugestivo. Una magnitud local de geometría queda determinada relacionalmente por la orientación y el retorno del estado completo. El área invariante restringe esa dependencia relacional y fija su consistencia.\par

La versión HMT del principio relacional consiste en algo más preciso que la máxima «la materia determina la geometría»: la orientación de la acción y la respuesta gravitatoria se codeterminan de tal manera que la geometría elemental permanezca invariante.\par

La acción contiene la historia. La gravedad ajusta la respuesta. El área conserva el acuerdo.\par

